\newcommand{\Val}{\mathrm{Val}}
\newcommand{\DVal}{\mathrm{DivVal}}

Given a variety $X$, a \emph{($\mathbb{Q}$-)valuation} is a map $v : \mathbb{C}(X)^\times \to \mathbb{Q}$ that satisfies, for all rational functions $f_1, f_2 \in \mathbb{C}(X)$,
\begin{itemize}
	\item $v(\mathbb{C}^\times) = 0$ (and as a convention, $v(0) = +\infty$);
	\item $v(f_1 + f_2) \geq \min(v(f_1), v(f_2))$; and
	\item $v(f_1 \cdot f_2) = v(f_1) + v(f_2)$.
\end{itemize}

We denote the \emph{center} of $v$ (if it exists) by $C_X(v) \subseteq X$. The center of $v$ is a closed irreducible subvariety of $X$ satisfying, for every rational function $f \in \mathbb{C}(X)$, $v(f) > 0$ iff $f = 0$ on $C_X(v)$. The uniqueness of center is guaranteed by separatedness of varieties.

A valuation $v$ is called a \emph{valuation on $X$} if the center exists. The trivial valuation on $X$ always exists, whose center is $X$ itself. If the variety $X$ is complete (e.g., projective), every valuation on $\mathbb{C}(X)$ is a valuation on $X$. We denote all the valuations on $X$ by $\Val_X$.

A Cartier divisor $D$ is locally defined by a single rational function, so we can define $v(D) \in \mathbb{Z}$; similarly, for a $\mathbb{Q}$-Cartier $\mathbb{Q}$-divisor $D$, we have $v(D) \in \mathbb{Q}$. If $D$ is effective, we have the inequality $v(D) \geq 0$. Valuations are $\mathbb{Q}_{\geq 0}$-linear on effective $\mathbb{Q}$-Cartier $\mathbb{Q}$-divisors.

Given a birational morphism $\pi: Y \to X$, a \emph{prime divisor $E$ over $X$} is defined as a prime divisor of $Y$. The prime divisor $E$ defines a valuation $\ord_E$ on $X$, whose center is the image $\pi(E)$. A \emph{divisorial ($\mathbb{Q}$-)valuation} $v$ is of the form $c \: \ord_E$ (for $c \in \mathbb{Q}_{\geq 0}$), and the space of divisorial valuations on $X$ is denoted by $\DVal_X$.

% equivariant settings
For a variety $X$ with an algebraic group $G$ acting on, a valuation $v$ is \emph{$G$-invariant} if it satisfies $v(g \cdot f) = v(f)$ for every $g \in G$ and $f \in \mathbb{C}(X)$. The space of $G$-invariant valuations is denoted by $\Val^G_X$, and similarly we have the space of $G$-invariant divisorial valuations $\DVal^G_X$.